\begin{abstract}
Public support for charging infrastructure must reconcile expanding coverage with the activity sustained by the available network. We examine this relationship using 2016--2023 city-week records from one anonymized operator network and locally applicable policy documents. Comparisons hold city membership fixed, distinguish positive activity from source-coded zeros, and relate city trajectories to facility-weighted network intensity and conditional annual and quarterly relationships. In the broad 2022--2023 comparison, activity grows by 15.32\% and observed facilities by 45.70\%, reducing network-wide intensity by 20.85\%. The 117 cities with coded-zero activity at both endpoints account for 52\% of the 225-city set but only 3.47\% of endpoint facilities. Growth types also change: among 52 matched cities, 17 of 27 early faster-growth cities subsequently record slower positive growth or activity contraction. Annual PPML estimates and the quarterly endpoint are positive and subproportional; the annual log-OLS interval includes unit elasticity ($p=0.260$), while the additional quarterly follow-through interval includes zero. Policy texts assign distinct information, maintenance and continuity responsibilities. Information/supervision is positively associated with recorded activity, whereas all operation/maintenance intervals include zero. These findings connect the assessment of infrastructure to both the scale and persistence of recorded service. Coverage provision, activity growth and operational continuity require separate evaluation criteria tied to the responsibilities public support actually specifies.
\end{abstract}
\noindent\textbf{Keywords:} electric vehicle charging; infrastructure expansion; activity intensity; distributed lags; operational policy
